%%%%%%%%%%%%%%%%%%%%%%%%%%%%%%%%%%%%%%%%%
% University Assignment Title Page 
% LaTeX Template
% Version 1.0 (27/12/12)
%
% This template has been downloaded from:
% http://www.LaTeXTemplates.com
%
% Original author:
% WikiBooks (http://en.wikibooks.org/wiki/LaTeX/Title_Creation)
%
% License:
% CC BY-NC-SA 3.0 (http://creativecommons.org/licenses/by-nc-sa/3.0/)
% 
%
%%%%%%%%%%%%%%%%%%%%%%%%%%%%%%%%%%%%%%%%%
%----------------------------------------------------------------------------------------
%	PACKAGES AND OTHER DOCUMENT CONFIGURATIONS
%----------------------------------------------------------------------------------------
\usepackage[a4paper,hmargin=2.5cm,vmargin=2.0cm,includeheadfoot]{geometry}
\usepackage{textpos}
\usepackage[numbers]{natbib} % for bibliography
\usepackage{tabularx,longtable,multirow,caption,subcaption}%hangcaption

\usepackage{url} % URLs
\usepackage[english]{babel}
\usepackage{amsmath}
\usepackage{graphicx}
\usepackage{dsfont}
\usepackage{epstopdf} % automatically replace .eps with .pdf in graphics
\usepackage{backref} % needed for citations
\usepackage{array}
\usepackage{latexsym}
\usepackage{wrapfig}
\usepackage[pdftex,pagebackref,hypertexnames=false,colorlinks]{hyperref} % provide links in pdf
\usepackage{parskip}
\usepackage[ruled,vlined]{algorithm2e}
\usepackage{enumitem}
\setlist[itemize]{left=0pt,labelsep=0.5em}

% Colored Quote
\usepackage{xcolor}
\usepackage{mdframed}
\newmdenv[
  backgroundcolor=gray!10, % light gray background
  linecolor=gray!50,       % border color
  skipabove=10pt,
  skipbelow=10pt,
  innertopmargin=6pt,
  innerbottommargin=6pt,
  innerleftmargin=10pt,
  innerrightmargin=10pt
]{colquote}

\hypersetup{pdftitle={},
  pdfsubject={}, 
  pdfauthor={},
  pdfkeywords={}, 
  pdfstartview=FitH,
  pdfpagemode={UseOutlines},% None, FullScreen, UseOutlines
  bookmarksnumbered=true, bookmarksopen=true, colorlinks,
    citecolor=black,%
    filecolor=black,%
    linkcolor=black,%
    urlcolor=black}

\usepackage[all]{hypcap}

% Table
\usepackage{tabularx} % For the tabularx environment and X column type
\usepackage{array}    % For advanced column specifications
% A new column type for centered, wrapping text in tabularx
\newcolumntype{C}{>{\centering\arraybackslash}X}

\usepackage{fancyhdr} % page layout
\fancypagestyle{plain}{%
    \fancyhf{} % clear all header and footer fields
    \fancyfoot[C]{\sffamily{\thepage}} % page number centered
}

%\usepackage{color}
%\usepackage[tight,ugly]{units}
%\usepackage{float}
%\usepackage{tcolorbox}
%\usepackage[colorinlistoftodos]{todonotes}
% \usepackage{ntheorem}
% \theoremstyle{break}
% \newtheorem{lemma}{Lemma}
% \newtheorem{theorem}{Theorem}
% \newtheorem{remark}{Remark}
% \newtheorem{definition}{Definition}
% \newtheorem{proof}{Proof}


%%% Default fonts
\renewcommand*{\rmdefault}{bch}
\renewcommand*{\ttdefault}{cmtt}



%%% Default settings (page layout)
\setlength{\parindent}{0em}  % indentation of paragraph

\setlength{\headheight}{14.5pt}
\pagestyle{fancy}
\renewcommand{\chaptermark}[1]{\markboth{\chaptername\ \thechapter.\ #1}{}} 

\fancyhf{} % clear all header and footer fields
\fancyfoot[C]{\sffamily{\thepage}} % page number centered
\fancyhead[LE,RO]{\slshape \rightmark} % section name on even (left) and odd (right) pages
\fancyhead[LO,RE]{\slshape \leftmark} % chapter name on odd (left) and even (right) pages
\renewcommand{\headrulewidth}{0.4pt}
\renewcommand{\footrulewidth}{0pt}
\captionsetup{margin=10pt,font=small,labelfont=bf}


%--- chapter heading

% \makeatletter
% \def\@makechapterhead#1{%
%   \vspace*{10\p@}%
%   {\parindent \z@ \raggedright \sffamily
%     \interlinepenalty\@M
%     \Huge\bfseries \chaptername\ \thechapter\space #1\par\nobreak
%     \vskip 10\p@ % Reduced spacing to save space
%   }}

%---chapter heading for \chapter*  
% \def\@makeschapterhead#1{%
%   \vspace*{10\p@}%
%   {\parindent \z@ \raggedright
%     \sffamily
%     \interlinepenalty\@M
%     \Huge \bfseries  #1\par\nobreak
%     \vskip 10\p@ % Also reduced here for consistency
%   }}
% \makeatother

\makeatletter
\def\@makechapterhead#1{%
  {\parindent \z@ \raggedright \sffamily
    \interlinepenalty\@M
    \Huge\bfseries \chaptername\ \thechapter\ --\ #1\par\nobreak
    \vskip 20\p@
  }}
\makeatother

% Disable forced page breaks at start of chapters
\usepackage{etoolbox}
\patchcmd{\chapter}{\if@openright\cleardoublepage\else\clearpage\fi}{}{}{}

\allowdisplaybreaks